\documentclass[11pt]{article}
\usepackage[T1]{fontenc}
\usepackage[utf8]{inputenc}
\usepackage{lmodern,amsmath,amssymb,graphicx,geometry,microtype,setspace,booktabs,caption}
\usepackage[hidelinks]{hyperref}
\usepackage{array}
\geometry{letterpaper,margin=1in}
\setlength{\parindent}{0pt}
\setlength{\parskip}{0.55em}
\setlength{\emergencystretch}{2em}
\captionsetup{labelfont=bf,font=small}
\renewcommand{\refname}{References}
\begin{document}
\begin{center}{\Large\bfseries S1 Appendix: Mechanisms and proofs}\\[0.5em]Yule Choi\end{center}
\onehalfspacing
\paragraph{Reading guide.} A1--A4 establish the stationary reduction and observation distinction. A5--A8 give finite biochemical realizations and explicit GlnD/GlnE compensation families. A9--A12 establish constructive rates, the closed-cascade bound and topology limits. A13--A14 derive calibrated reporter intervals; A15 distinguishes these contributions from prior theory. Data provenance, numerical sensitivity and assay precision are in S2 Appendix. Equation identifiers inherited from the original derivations are retained to support comparison with the repository archive; section identifiers are specific to this appendix.
\section*{A1 Reaction network and the steady-state invariant}

Let E be free enzyme, L free effector, and $S_i$ a free target with i modified sites, for i=0,\ldots{},n. The core network is


\[
E+L \mathrel{\mathop{\rightleftharpoons}^{a_1}_{a_2}} EL, \tag{S1}
\]

\[
E+S_i \mathrel{\mathop{\rightleftharpoons}^{b_{1i}}_{b_{2i}}} B_i
\mathrel{\mathop{\longrightarrow}^{b_{3i}}} E+S_{i+1}, \tag{S2}
\]

\[
EL+S_{i+1} \mathrel{\mathop{\rightleftharpoons}^{g_{1i}}_{g_{2i}}} G_i
\mathrel{\mathop{\longrightarrow}^{g_{3i}}} EL+S_i,\qquad i=0,\ldots,n-1. \tag{S3}
\]

All rate constants are positive. The catalytic mechanism is distributive and each complex contains one enzyme and one target. Lower-case symbols denote the concentrations of the corresponding species, with $\varepsilon$=[EL], $\beta_i$=[$B_i$], and $\gamma_i$=[$G_i$]. Define


\[
A=\frac{a_1}{a_2},\quad
\widehat B_i=\frac{b_{1i}}{b_{2i}+b_{3i}},\quad
\widehat G_i=\frac{g_{1i}}{g_{2i}+g_{3i}},\quad
\rho_i=\frac{b_{3i}}{g_{3i}}. \tag{S4}
\]

The conserved totals are


\[
\begin{aligned}
T_E=e+\varepsilon+\sum_i(\beta_i+\gamma_i),\\
T_L=l+\varepsilon+\sum_i\gamma_i,\\
T_S=\sum_{i=0}^{n}s_i+\sum_i(\beta_i+\gamma_i).
\end{aligned} \tag{S5}
\]

At a positive steady state, the $G_i$ balances give $\gamma_i$=$\widehat G_i$ $\varepsilon$$s_{i+1}$. Substituting their balances into the EL balance cancels all target-binding contributions, leaving $a_1$el=$a_2$$\varepsilon$. Thus $\varepsilon$=Ael. The corresponding $B_i$ balances give $\beta_i$=$\widehat B_i$ e$s_i$.


The target modification states form a chain. Summing the balances on one side of each cut gives equality of the forward and reverse catalytic fluxes across it: $b_{3i}$$\beta_i$=$g_{3i}$$\gamma_i$. Substitution therefore yields


\[
\frac{s_{i+1}}{s_i}=\frac{K_{i+1}}{l},\qquad
K_{i+1}=\frac{b_{3i}\widehat B_i}{A g_{3i}\widehat G_i}. \tag{S6}
\]

This relation does not assume fast binding or small enzyme concentration. It is an identity at each positive steady state of the stated network. It does not establish how many such states exist, whether they are stable, or how fast they are approached. The enzyme pool cancels from ratios; absolute free-target concentrations can still vary with conserved totals.


\section*{A2 Distribution and local factorization}

Suppose more generally that adjacent ratios have the separable form $s_i$/$s_{i-1}$=$K_i$/f(l), with $K_i$ independent of input. This is an assumption about the effective regulation in an extended network, not a consequence of bifunctionality alone. Fix a consistent concentration unit so that $K_i$/f(l) is dimensionless. Let u=1/f(l), $c_0$=1, and $c_i$ be the product of $K_1$ through $K_i$. Then


\[
s_i=s_0c_i u^i,\qquad P(u)=\sum_{i=0}^{n}c_i u^i,\qquad
p_i=\frac{c_i u^i}{P(u)}. \tag{S7}
\]

Let $\mu$=$\Sigma$\_i i $p_i$ and $\theta$=$\mu$/n. Direct differentiation gives d $p_i$/dlog u=(i-$\mu$)$p_i$, and hence d$\mu$/dlog u=Var(i). Consequently,


\[
-\frac{d\log[\theta/(1-\theta)]}{d\log l}
=\frac{n\operatorname{Var}(i)}{\mu(n-\mu)}\frac{d\log f}{d\log l}
=\nu\eta. \tag{S8}
\]

The product applies to the physical free-target mean at a common input. It does not imply a product rule for fitted exponents or response-range coefficients. If f is increasing, the mean decreases and the chosen negative derivative is positive. An increasing readout uses the corresponding positive derivative.


For modification count i in [0,n], the inequality $i^2$$\le$ni implies Var(i)$\le$$\mu$(n-$\mu$). Therefore 0<$\nu$$\le$n. With n>1 and strictly positive weights on intermediate states, the upper inequality is strict. Its supremum is approached as intermediate weights vanish and the distribution becomes concentrated on i=0 and i=n. For n=1, $\nu$=1 identically.


For identical independent sites, $K_i$=(n-i+1)$\widehat K$/i. Therefore P(u)=$(1+\widehat K u)^n$ and $\theta$=$\widehat K$ u/(1+$\widehat K$ u). Its ladder factor is one for every u. Conversely, if $\nu$=1 throughout an open input interval, then dlog[$\theta$/(1-$\theta$)]/dlog u=1. Integration gives $\theta$=Ku/(1+Ku) for a positive constant K, and integrating uP'/P=n$\theta$ gives P=$(1+Ku)^n$ after imposing P(0)=1. Equality $\nu$=1 at a single input does not establish this global form.


The state-distribution example in the archived state-distribution figure (legacy S8 Fig D) uses n=12 and weights proportional to binomial(12,i) exp[J$(i-6)^2$]. At u=1, symmetry fixes $\theta$=1/2 for every J. Choosing J$\approx$0.109876 gives $\nu$=2.2, whereas J=0 gives $\nu$=1. Both distributions are explicitly normalized in S1 Data. This is an illustrative family rather than an inferred biochemical ladder.


\section*{A3 Opposing effectors and inactive complexes}

For two free effectors L\_1 and L\_2 binding E independently into mutually exclusive competent forms, $\varepsilon$\_1=A\_1el\_1 and $\varepsilon$\_2=A\_2el\_2. If EL\_1 catalyzes the forward reaction and EL\_2 the reverse, flux balance gives


\[
\frac{s_{i+1}}{s_i}=\frac{b_{3i}\widehat B_i A_1}{g_{3i}\widehat G_i A_2}\frac{l_1}{l_2}. \tag{S9}
\]

Thus f=l\_2/l\_1 after absorbing constant prefactors into the ladder. If both effectors depend on an external signal, the switch factor is the log derivative of their ratio with respect to that signal. Additional direct regulation can be absorbed into a common f only if the resulting factor is shared across all ladder steps. Productive target-bridging complexes or modification-state-dependent regulatory functions can break this separation.


A reversible dead-end species whose only reactions are formation and dissociation has zero net formation flux at steady state. Adding finitely many such species therefore leaves the balance identities underlying Eq. (S6) unchanged at fixed free inputs. This statement concerns the invariant, not the free-to-total map. Added complexes change conservation relations, and complexes carrying more than one target require a corresponding change to target-sequestration bounds. A reduced-rate catalytic complex is not a dead end and is not covered by this argument.


\section*{A4 Free and total observables}

Assume that no enzyme molecule can carry more than h effector molecules, including all complexes being counted. Writing $T_L$=l+L\_bound gives 0$\le$L\_bound$\le$h$T_E$ and


\[
\max(0,T_L-hT_E)\le l\le T_L. \tag{S10}
\]

For the core one-target-per-complex model, let q be the fraction of target held in enzyme complexes and $\theta_b$ its mean fractional modification. Then $\theta_{\mathrm{total}}$=(1-q)$\theta_{\mathrm{free}}$+q$\theta_b$. Since both fractions lie in [0,1],


\[
|\theta_{\rm total}-\theta_{\rm free}|\le q\le
\min(1,T_E/T_S). \tag{S11}
\]

Neither concentration nor amplitude bounds alone control derivatives. A finite-range statement can nevertheless be obtained without bounding a derivative. On a monotone branch, let $l_a$<$l_b$ be the free-input values at two prescribed free-target response thresholds. Put $\delta$=h$T_E$. Their corresponding total inputs satisfy $T_a$ in [$l_a$,$l_a$+$\delta$] and $T_b$ in [$l_b$,$l_b$+$\delta$]. Therefore


\[
\frac{l_b}{l_a+\delta}\le \frac{T_b}{T_a}\le\frac{l_b+\delta}{l_a}. \tag{S12}
\]

If $l_b$>$l_a$+$\delta$, the induced response-range coefficient obeys


\[
\frac{\log81}{\log[(l_b+\delta)/l_a]}
\le n_{H,\rm total\ input}^{\rm range}
\le\frac{\log81}{\log[l_b/(l_a+\delta)]}. \tag{S13}
\]

The thresholds here belong to the same free-target response; changing the ordinate to total target requires a separate treatment. This qualification avoids using a small sequestration pool as a universal local-slope bound. It also separates enzyme dilution from fast binding, which is a dynamical assumption absent from Eq. (S10).


\section*{A5 Finite ligand and catalytic states realizing the population ambiguity}

All quantities in this section refer to free target and buffered free glutamine unless otherwise stated. Three enzyme states have binding weights $w=(1,G/K_1,G^2/(K_1K_2))$ and nonnegative forward and reverse catalytic specificities $\alpha_j,\beta_j$. Their effective forward/reverse ratio is $A/B$, where $A=\sum_j\alpha_jw_j$ and $B=\sum_j\beta_jw_j$. Independent modification sites yield $\theta=A/(A+B)$ and binomial free-target state fractions.

One explicit finite realization uses $E_0+G\rightleftharpoons E_1$ and $E_1+G\rightleftharpoons E_2$, with ligand binding only to the free enzyme states. For each $j=0,1,2$ and $i=0,1,2$, include
\[
E_j+S_i\rightleftharpoons B_{ji}\longrightarrow E_j+S_{i+1},\qquad
E_j+S_{i+1}\rightleftharpoons C_{ji}\longrightarrow E_j+S_i.
\]
Choose forward specificities $(3-i)\alpha_j$ and reverse specificities $(i+1)\beta_j$. At steady state, eliminating each complex cancels its contribution to the free-enzyme balance. The ligand-binding ratios therefore retain weights $w$, and the target fluxes give the binomial ladder with input $A/B$. For example, in consistent nondimensional units, $k_{\rm off}=k_{\rm cat}=1$ and $k_{\rm on}=2$ times the desired specificity provide such a realization. Buffered nucleotide cofactors are absorbed into catalytic rates. This establishes the free-target result with explicit complexes; it does not identify closed-system total modification with free modification. The supplied verification constructs ten positive/nonnegative steady states and checks every species balance. Conserved totals are computed for each state and are not claimed to remain fixed across these constructed glutamine inputs.

Let $r\le\min_j \alpha_j/(\alpha_j+\beta_j)$ and set $\alpha_j^{\rm R}=(1-r)\alpha_j-r\beta_j$ and $\beta_j^{\rm R}=\beta_j$. Then $A_{\rm R}=(1-r)A-rB$ and $A_{\rm R}+B=(1-r)(A+B)$, so $q=A_{\rm R}/(A_{\rm R}+B)=(\theta-r)/(1-r)$. A separate fully modified species with free fraction $r$ and no GlnD-mediated turnover gives the same total free mean. The refractory species is assumed available to the stated downstream recognition law. It is neither an experimentally established species nor a dynamically exchanging state. Vanishing catalytic coefficients denote omitted reactions. For $r=\alpha_2/(\alpha_2+\beta_2)$, the accessible model has $\alpha_2^{\rm R}=0$.

For the independent-binding implementation, set $K_1=K/2$, $K_2=2K$, $\alpha_1=\alpha_2$ and $\beta_1=\beta_0$. Thus
\[
U(G)=\alpha_2+\frac{\alpha_0-\alpha_2}{(1+G/K)^2},\qquad
R(G)=\beta_0+(\beta_2-\beta_0)\left(\frac{G/K}{1+G/K}\right)^2.
\]
These capacities are monotone for $\alpha_0\ge\alpha_2$ and $\beta_2\ge\beta_0$. The refractory transformation preserves monotonicity because it subtracts an increasing reverse capacity from a decreasing forward capacity. The model has independent ligand binding with state-dependent catalytic gating; it does not claim that enzyme activities respond independently to each binding event.

The fit parameters are $q_0$, $q_2$, and $K$, where $q_j=\alpha_j/(\alpha_j+\beta_j)$. Fix the arbitrary common scale by $\alpha_0+\beta_0=1$. With transported specificity fold $F=6.752484$, set $d_2=F(1-q_0)/(1-q_2)$, $\alpha_2=d_2q_2$, and $\beta_2=d_2(1-q_2)$. Bounds are $0.9\le q_0\le0.999999$, $10^{-6}\le q_2\le0.499999$, and $0.001\le K\le10$ mM. Three initializations are supplied. Every leave-one-coordinate-out fit is independently optimized. Outputs report all fitted coefficients and errors. No microscopic binding data were included in the objective.

For comparison, positive rational means are parameterized as
\[
\theta(G)=\frac{q_0+q_1G/(s\sqrt C)+q_2(G/s)^2}{1+G/(s\sqrt C)+(G/s)^2}.
\]
Here $C$ is an effective coefficient ratio, not generally $K_1/K_2$. A profile restricts $C$ to upper bounds 0.25, 1, 10, 100, and 10000 and fits ordered $q_0\ge q_1\ge q_2$. The upper boundary is reached in every scenario; this is evidence of limited identification within that chosen family, not an estimate of binding cooperativity. For a second, fixed-$C=100$ construction, set $q_1=q_0$, fix $F$ as above, and impose activity half-range $h=0.080$ mM using
\[
d_2=\frac{h^2}{s^2+sh/\sqrt C},\qquad q_0=1-d_2(1-q_2)/F.
\]
Only $q_2$ and $s$ are fitted. Choose $d_1=1$; the actual binding constants are $K_1=s\sqrt C$ and $K_2=sd_2/\sqrt C$. Their ratio is approximately 4500 in this transferred-kinetics construction. This is a strong, unmeasured cooperativity requirement. The independent-binding fit demonstrates that such a requirement is not inherent to the population equivalence.

\section*{A6 GlnD occupancy--catalysis transformations}

Use $E_0,E_1,E_2$ with the finite enzyme--target-complex realization in A5. Let $D_1=1+2z+z^2$, $z=G/K$, and $A=a_0+2a_1z+a_2z^2$, $B=b_0+2b_1z+b_2z^2$. The reference sets $a_1=a_2$ and $b_1=b_0$. For $\lambda>0$, change the middle binding weight to $2z/\lambda$ and middle catalytic coefficients to $\lambda a_1,\lambda b_1$. Equivalently, $K_1=\lambda K/2$ and $K_2=2K/\lambda$. Then $A$ and $B$ are unchanged, while $D_\lambda=1+2z/\lambda+z^2$. Both reduced capacities acquire the same positive multiplier
\[
\frac{U_\lambda}{U_1}=\frac{R_\lambda}{R_1}
=\frac{D_1}{D_\lambda}.
\]
The free-target ratio, mean, and all binomial modification-state fractions are unchanged. The finite-complex stationary identity concerns free targets. The displayed kinetic multiplier additionally uses a fast-binding, unsaturated-conversion reduction in which the regulated free-enzyme ensemble provides the common rate scale. It should not be applied unchanged to a strongly sequestered closed-total assay.

Ordered state activities ensure monotone capacities. For the fitted reference, the interval $1\le\lambda\le\min(b_2/b_0,a_0/a_1)$ preserves $a_0\ge\lambda a_1\ge a_2$ and $b_0\le\lambda b_1\le b_2$. Here the upper bound is the transported UR-fold scenario 6.752484, so $\lambda=1,2,4$ are all admissible. These are constructed alternatives, not three separately optimized explanations. The reference is refitted to the same twelve coordinates using three starts, with $q_0\in[0.9,0.999999]$, $q_2\in[10^{-6},0.499999]$, and $K\in[0.001,10]$ mM. The reference catalytic scale is arbitrary; only the ratios enter the mean fit. Complete coefficients are reported in S1 Data.

At $G=K$, the singly liganded enzyme fraction is $1/(1+\lambda)$ and the capacity multiplier is $2\lambda/(1+\lambda)$. Thus the $\lambda=1$ and 4 models have fractions 0.5 and 0.2 and a speed ratio of 1.6. Both have mean ligand occupancy exactly one at that input. At $G=K/3$, their mean ligand occupancies are 0.5 and approximately 0.304; at $G=3K$, they are 1.5 and approximately 1.696. These flanking inputs are useful for an average-binding assay, whereas a central single-state or rate measurement addresses the same ambiguity differently. The fitted $K$ is a conditional model scale, not a directly measured glutamine dissociation constant.

For the step-response illustration, $\dot\theta=s_D[U(G)(1-\theta)-R(G)\theta]$ and $\dot Y=s_E[v_{\rm AT}(1-Y)-v_{\rm AR}Y]$. The initial state is the steady state at $G=0.5$ mM, and the final free glutamine is 0.05 mM. $T=5\,\mu$M, and the downstream recognition law is $P=T(1-\theta)^3$, $U=T\theta$. Shared scales $s_D,s_E$ make the reference relaxation rates at the final steady state unity; transformed models retain those same scales. This is a prediction in nondimensional reference units, not a fit to experimental time courses. The source functions and generated file include both the PII and GS trajectories.

\section*{A7 Identifiable combinations in the six-state GlnE model}

The state activities and topology follow Fig 12 of Jiang, Mayo and Ninfa (2007). The ordered state weights are $(1,g,p,u,pg/\alpha_1,pu/\alpha_2)$ for $E$, $E$--Gln, $E$--PII, $E$--PII-UMP, $E$--Gln--PII, and $E$--PII--PII-UMP. AT activities are $k_P(0,\beta,1,0,\beta',0)$, and AR activities are $k_R(0,0,0,1,0,0)$. Glutamine and PII-UMP do not occupy the enzyme simultaneously in this topology. The recognition terms $P,U$ can be externally controlled free inputs. In the coupled illustrations, they are assigned the specified unmodified-trimer and count-weighted reverse inputs. This recognition assignment is distinct from the source topology itself.

The common polynomial cancels from $v=v_{\rm AT}/v_{\rm AR}$. Accordingly, $\alpha_2$ is absent from this steady drive and $\alpha_1,\beta'$ enter only through $\beta'/\alpha_1$. For $c>0$, the transformation $(\alpha_1,\beta')\mapsto(c\alpha_1,c\beta')$ multiplies both catalytic rates by $Z_E/Z'_E$ but leaves their ratio unchanged. Changing $\alpha_2$ has the same common-rate effect with a different occupancy dependence. These identities preserve any downstream steady-state construction depending only on that drive, not only the independent-site GS example. Kinetic predictions require the conversion assumptions specified in the main Methods. Summing the three states bound to unmodified PII gives a binding isotherm with halfpoint
\[
K_{P,\mathrm{eff}}=K_P\frac{1+g+u}{1+g/\alpha_1+u/\alpha_2}.
\]
This prediction depends on $\alpha_1$ separately from $\beta'$, providing a complementary measurement. It is a predicted binding halfpoint, not an assumption that every published apparent activation constant equals a microscopic dissociation constant.

A six-state continuous-time generator independently verifies the binding probabilities. Choose edges $E\leftrightarrow E$--Gln, $E\leftrightarrow E$--PII, $E\leftrightarrow E$--PII-UMP, $E$--Gln$\leftrightarrow E$--Gln--PII, $E$--PII$\leftrightarrow E$--Gln--PII, $E$--PII$\leftrightarrow E$--PII--PII-UMP, and $E$--PII-UMP$\leftrightarrow E$--PII--PII-UMP. Unit reverse rates and forward rates $(g,p,u,p/\alpha_1,g/\alpha_1,u/\alpha_2,p/\alpha_2)$ satisfy detailed balance with the stated weights. Stationary probabilities from linear solves agree with normalized weights in 250 random positive cases. These generator rates verify the equilibrium topology; they are not inferred binding-time constants.

Parameter provenance is kept separate from calculation. The source's A1/A2 fitting example reports $K_G=15.6$ mM, $K_P=0.18\,\mu$M and $\alpha_1=0.17$. The selected $K_U=0.35\,\mu$M is an apparent activation-scale scenario from Table 3. Values $\beta=1$, $\beta'=10$, and $\alpha_2=4$ are chosen illustrations. The source reports $\alpha_2$ estimates 2.17--8.85 from other global fits; these motivate variants but are not a shared-condition confidence interval. The reference $k_P/k_R$ is set once so that the independent-GS mean equals 0.5 at $G=0.53$ mM and $T=5\,\mu$M. No other midpoint or response coefficient is fitted. The transformations retain this gain. Thus the steady-curve identities are exact mathematical consequences, while the numerical occupancy and time-scale differences are conditional predictions.

For fixed $G,\theta$ and the coupled recognition law, write
\[
v(T)=\frac{k_PK_U}{k_RK_P\theta}(1-\theta)^3
\left(1+\frac{\beta'G}{\alpha_1K_G}\right)
+\frac{k_PK_U\beta G}{k_RK_G\theta}\frac1T.
\]
Only the PII-independent AT state supplies the inverse-$T$ term. The remaining pathway terms are concentration invariant at fixed free-state fractions. With independent GS sites, $v=Y/(1-Y)$; for a different downstream observation law the affine prediction applies to the effective drive rather than automatically to measured GS odds. In S5 Fig D, the analytic $T\to\infty$ intercept is subtracted only to display the concentration-dependent contribution. The three historical cascade assays do not by themselves test this relation because free inputs were not held fixed and GlnD concentrations also changed.

Let the normalized positive contributions to the AT numerator be $w_P,w_G,w_{PG}$ from $p,\beta g,\beta'pg/\alpha_1$. At fixed $P,U$,
\[
\left.\frac{\partial\log v}{\partial\log G}\right|_{P,U}
=w_G+w_{PG}\in[0,1].
\]
Along the coupled input path, define $a=-d\log\theta/d\log G$ and $b=3\theta a/(1-\theta)$. Then
\[
\frac{d\log v}{d\log G}=a+(w_P+w_{PG})b+(w_G+w_{PG}).
\]
The first two terms are mediated by PII and the last is the direct glutamine contribution. A centered logarithmic finite difference checks this identity over 701 inputs. This is a local drive decomposition; endpoint-normalized response-range coefficients and nonindependent GS ladder factors require separate calculations.

The GlnD transformation and the two GlnE transformations combine without changing any steady target output under these assumptions. Thirty-six combinations of target total and mechanism are verified numerically. They define complementary experiments: resolve GlnD occupancy or PII relaxation to test the upstream family; measure GlnE binding and AT/AR initial rates across controlled ligand inputs to separate affinity and catalytic synergy; and vary free target total while monitoring its state fractions to test the PII-independent route. Parameter constraints and assay conservation must accompany a physiological interpretation.


\clearpage

\clearpage
\section*{A8 Finite reaction realization, total-readout bound, and catalytic coupling}

Let $X_0,X_1$ be two free states of one modification unit. Three free enzyme states exchange glutamine with finite rates. For $j=0,1,2$ introduce
\[
 E_j+X_0\mathop{\rightleftharpoons}^{k_{\rm on,j}^+}_{k_{\rm off,j}^+} C_j^+
 \mathop{\longrightarrow}^{k_{\rm cat,j}^+} E_j+X_1,
\qquad
 E_j+X_1\mathop{\rightleftharpoons}^{k_{\rm on,j}^-}_{k_{\rm off,j}^-} C_j^-
 \mathop{\longrightarrow}^{k_{\rm cat,j}^-} E_j+X_0.
\]
All inputs are buffered free glutamine; enzyme and target totals are conserved. Set $h_j^\pm=k_{\rm on,j}^\pm/(k_{\rm off,j}^\pm+k_{\rm cat,j}^\pm)$, $a_j=h_j^+k_{\rm cat,j}^+$, and $b_j=h_j^-k_{\rm cat,j}^-$. With ligand exchange restricted to free enzyme, the steady complex equations give $C_j^+=h_j^+E_jX_0$ and $C_j^-=h_j^-E_jX_1$. Each complex's net association minus release-and-catalysis is zero, leaving the free-enzyme equations $Q_EE=0$. Any positive common multiplier of the ligand-binding generator changes its relaxation but not its stationary weights $w=(1,G/K_1,G^2/(K_1K_2))$. Consequently
\[
 \frac{X_1}{X_0}=\frac{\sum_j a_jw_j}{\sum_j b_jw_j}
\]
is exact for this full reaction system, without a rapid-equilibrium or unsaturated-target approximation. The earlier $\lambda$ transformation preserves it. This statement concerns the stationary equations, not uniqueness for every possible topology or a universal first-order time course.

Let $q=X_1/(X_0+X_1)$, $x=X_0+X_1$, and let $o_j=w_j/\sum w$. Define $H_+=\sum o_jh_j^+$, $H_-=\sum o_jh_j^-$, and $H=(1-q)H_++qH_-$. The free enzyme pool is $E_f=E_T/(1+Hx)$ and target conservation gives
\[
 Hx^2+[1+H(E_T-T_T)]x-T_T=0.
\]
Its positive root yields all 11 concentrations. The measured total modified fraction is $Y_T=(qx+E_fqxH_-)/T_T$. If $H_+=H_-$, then $Y_T=q$ exactly, including when most enzyme is complexed. Equality $h_j^+=h_j^-$ state by state is sufficient.

For general affinities write $B=\sum(C_j^++C_j^-)$ and $M_B=\sum C_j^-$. Since $0\leq M_B\leq B\leq\min(E_T,T_T)$,
\[
Y_T=q+\frac{M_B-qB}{T_T},\qquad
q-q\epsilon\leq Y_T\leq q+(1-q)\epsilon,
\quad \epsilon=\min(1,E_T/T_T).
\]
Two models with the same $q,E_T,T_T$ lie in the same interval of width $\epsilon$, proving the pairwise bound in the main text. If one enzyme can bind more than one counted modification unit, the upper bound on $B$ must be adjusted to that stoichiometry. In particular, a monovalent unit calculation must not be relabelled a fully resolved trimer sequestration model. No derivative bound follows from the concentration bound alone.

The numerical realization sets $k_{\rm cat}^+=a/h^+$, $k_{\rm cat}^-=b/h^-$, $k_{\rm off}=k_{\rm cat}$, and $k_{\rm on}=2h k_{\rm cat}$. A shared scale makes the reference summed capacity at $G=K$ equal to one. The illustrative asymmetric affinities are $h^+=(0.2,2,10)$ and $h^-=(5,0.3,0.1)$ in reciprocal target-concentration units. Target total is one; enzyme/target ratios range from $10^{-4}$ to one. Random verification uses 300 positive cases with all affinity coordinates independently spanning $10^{-2}$--$10^2$, binding-rate multipliers $10^{-3}$--$10^3$, and target totals $10^{-1}$--$10$. Full balances, conserved totals, the pairwise bound, and the symmetric-affinity identity are checked. Two BDF integrations from initially free enzyme and unmodified target converge to the independently constructed steady states.

For the finite-binding transient in S6 Fig C, enzyme fractions obey $\dot o=Q_Eo$ and the dilute target obeys $\dot q=(o\cdot a)(1-q)-(o\cdot b)q$. The glutamine step is $10K\to K$, initialized at the preceding steady state. Binding multipliers span 0.01--100. These equations are an unsaturated mean-field conversion model with explicit regulatory relaxation, distinct from the 11-concentration saturation model. The units are reference times, not fitted seconds or minutes.

\textbf{Exchange in substrate-bound complexes.} Adding ligand exchange to $C^+$ and $C^-$ retains the constructed steady state if their generators annihilate the vectors $w_jh_j^+$ and $w_jh_j^-$. This is a sufficient kinetic compatibility condition and can hold for nonzero exchange rates. Thermodynamic compatibility of binding alone is weaker. Write $\rho_j=k_{\rm off,j}/k_{\rm cat,j}$; the equilibrium association constant is $k_{\rm on}/k_{\rm off}=h_j(1+\rho_j)/\rho_j$, which need not be proportional to the catalytic-complex weight $h_j$ across states.

To test this distinction without arbitrary inconsistent binding cycles, retain the same $h,a,b$ but choose $\rho^+=(0.1,1,10)$ and $\rho^-=(10,1,0.1)$. Set $k_{\rm on}=(1+\rho)h k_{\rm cat}$. For each complex chain, take reverse ligand-exchange rate $\eta$ and forward rate equal to $\eta$ times the free-chain association ratio times the adjacent ratio of $h(1+\rho)/\rho$. All noncatalytic binding cycles then satisfy the equilibrium ratio constraints. Catalysis can nevertheless sustain regulatory currents; the buffered nucleotide-driven conversion is not assumed at thermodynamic equilibrium.

Positive log-concentration stationary solves cover $G/K=0.3,1,3,10,30$, $\eta=0$ and 25 logarithmically spaced values from $10^{-4}$ to $10^2$, $E_T/T_T=0.1$, and both $\lambda$ values. The full ODE residual and both conservation laws are checked; independent BDF integrations at $G=10K,\eta=1$ verify two stationary solutions. In this scenario the maximum scanned free-fraction gaps are about 0.035 at $10K$ and 0.070 at $30K$, occurring near $\eta=0.1$. These values establish a conditional counterexample to unrestricted topology invariance; they are not estimates of biological rates or evidence that this coupling occurs in GlnD.

\clearpage
\section*{A9 Constructive stationary compensation and admissibility}

Consider a connected free-enzyme regulatory network with a positive stationary weight vector $w(x)$ at fixed free ligands $x$. Binding may be reversible with thermodynamically consistent cycle products; more generally the free-state linear network must admit these weights. A catalytic intermediate $C_{ji}^{+}$ is formed from regulatory state $E_j$ and target $S_i$ and dissociates or catalyzes to $E_j+S_{i+1}$. All catalytic branches return the same regulatory state. Regulatory exchange on a catalytic intermediate is absent. At any positive stationary point,
\[
 C_{ji}^{+}=h_{ji}^{+}E_jS_i,\qquad
 h_{ji}^{+}=\frac{k_{\mathrm{on},ji}^{+}}{k_{\mathrm{off},ji}^{+}+k_{\mathrm{cat},ji}^{+}},\qquad
 a_{ji}^{+}=h_{ji}^{+}k_{\mathrm{cat},ji}^{+},
\]
with analogous reverse branches. Eliminating these branches from the free-enzyme balance leaves its regulatory network, so $E_j=e_0w_j$. Zero target flux across each ladder cut gives
\[
 \frac{S_{i+1}}{S_i}=\frac{\sum_jw_ja_{ji}^{+}}{\sum_jw_ja_{j,i+1}^{-}}.
\]
For constants $s_j>0$, the transformation $w'_j=s_jw_j$, $(a_{ji}^{\pm})'=a_{ji}^{\pm}/s_j$ preserves every edge ratio. This holds at stationary state without assuming fast catalytic binding. Positive specificities and association coefficients have a finite realization: choose $k_{\rm cat}=a/h$, $k_{\rm off}=a/h$, and $k_{\rm on}=2a$. The binding constants generating $w'$ must be physically admissible within the chosen regulatory network; arbitrary independently rescaled equilibrium edges are not admissible if they violate a binding cycle. Changing the weights of states with no catalytic branches also leaves the flux numerators unchanged.

This is a sufficient construction, not a converse. More general equivalences can involve redistribution of states, distinct monomials, or non-scaling transformations. Moreover the free-enzyme normalization, protein sequestration, and relaxation times can change. The statement concerns normalized free-target stationary distributions; it does not establish dynamic unidentifiability or equality at fixed total ligand. Applying it to both cascade layers requires additional closure arguments because one layer's free protein pool is the next layer's ligand supply.

\section*{A10 Explicit state-resolved, protein-conserving cascade}

\textbf{Species and conserved quantities.} The buffered-input reference contains free PII trimers $P_i$, $i=0,\ldots,3$ (four species); free GS dodecamers $S_k$, $k=0,\ldots,12$ (thirteen); free GlnD states $D_j$, $j=0,1,2$ (three); eighteen GlnD--PII catalytic intermediates; ten GlnE regulatory states; and seventy-two GlnE--GS catalytic intermediates. These total 120 species and 300 directed reactions. The GlnE states are
\[
 E,\ EG,\ EP_0,\ EU_1,\ EU_2,\ EU_3,\ EGP_0,\ EP_0U_1,\ EP_0U_2,\ EP_0U_3.
\]
Here $U_i$ denotes a bound PII trimer carrying $i$ UMP groups, not a distinct free protein species. The EPU states contain two regulatory PII trimers. A GS-bound GlnE retains its regulatory ligands. The stoichiometric matrix has four exact conserved protein totals: PII, GS, GlnD, and GlnE. Modification counts in catalytic complexes refer to the pre-catalytic target state. Glutamine is buffered in the base construction; ATP/ADP, UTP and products are not explicit species. Thus ``protein-conserving'' does not mean metabolically or energetically closed.

\textbf{Regulatory binding.} GlnD has sequential binding weights $(1,G/K_1,G^2/(K_1K_2))$. The GlnE weights are
\[
 (1,g,p,u_1,u_2,u_3,pg/\alpha_1,pu_1/\alpha_2,pu_2/\alpha_2,pu_3/\alpha_2),
\]
where $g=G/K_G$, $p=P_0/K_P$, and $u_i=(i/3)P_i/K_U$. Summing the $u_i$ recovers the earlier six-class denominator. The recognition weight $i/3$ is a modeling choice, not a measured affinity of every partially uridylylated trimer. All four edges around each binding square are included with compatible equilibrium products. Reverse regulatory rates are one in arbitrary time units; forward rates are the corresponding ligand concentration divided by its dissociation constant. Regulatory exchange occurs only when the enzyme is free of its catalytic target. GlnD does not modify GlnE-bound PII in this construction.

\textbf{Catalytic edges.} GlnD forward and reverse specificities are $a_j(3-i)$ and $b_ji$, respectively. Base catalytic association coefficients in $\mu$M$^{-1}$ are $(0.2,2,10)$ forward and $(5,0.3,0.1)$ reverse, multiplied by the same available-site counts. GlnE ATase-active states are EG, EP, EGP, with specificity multipliers $k\beta$, $k$, and $k\beta'$; AR-active states are $EU_i$, each with multiplier one. Available GS sites give factors $12-k$ forward and $k$ reverse. For the interaction scenario, the forward factor is instead $(12-k)\exp[J(2k+1-12)]$, giving
\[
 c_k={12\choose k}\exp[J(k^2-12k)],\qquad \pi_k\propto c_kr^k.
\]
The GS catalytic association coefficient is 0.1 $\mu$M$^{-1}$ times the corresponding site factor. Each positive branch is realized using the rates in A9. With the direct route removed, EG catalytic branches and their twelve complexes are absent (the numerical state vector retains zero-valued slots). All other baseline rates are positive. For the productive-route extension, EPU states acquire ATase specificity $k\eta$; thirty-six further GS complexes give 156 species for $\eta>0$.

\textbf{Parameters and provenance.} Concentrations in the new code are uniformly micromolar. The transported GlnD reference is $K_1=26.02223$, $K_2=104.08892$, with specificity vectors
\[\begin{aligned}
 a&=(0.99889564,0.0012739725,0.0012739725),\\
 b&=(0.0011043614,0.0011043614,0.0074571831).
\end{aligned}\]
 Its $\lambda$ transformation multiplies $K_1,a_1,b_1$ by $\lambda$ and divides $K_2$ by $\lambda$. GlnE uses $K_G=15600$, $K_P=0.18$, $K_U=0.35$, $\alpha_1=0.17$, $\alpha_2=4$, $\beta=1$, $\beta'=10$, and $k=1.4877648555$, with the source/calibration distinctions stated in A6--A7 and the regulatory-state parameter records. The $\sigma$ transformation multiplies both $\alpha_1$ and $\beta'$ by $\sigma$. Catalytic association coefficients, binding time scales, recognition weights, protein-pool scenarios, and $J$ are chosen stress-test parameters, not newly estimated biological constants. Setting $\beta=0$ removes direct glutamine catalysis.

\textbf{Steady-state solution and verification.} At fixed free $G$, normalized free PII follows the binomial distribution with fraction $q_D=(w\cdot a)/(w\cdot(a+b))$. For candidate free pools $T$ and $V$, set $P_i=T\pi_i$ and $S_k=V\pi_k(r)$. Free-enzyme normalizations follow directly from their totals and the sum of regulatory and catalytic-complex weights. The two remaining equations impose total PII and GS conservation. We solve these equations numerically, then insert every constructed species into the independently enumerated mass-action vector field. No quasi-steady-state approximation is used in this stationary verification. The reaction generator records all species, stoichiometry, and rates and checks the conservation matrix exactly.

The paired sweep uses 82 distinct glutamine doses: 41 logarithmically spaced from 20 to 10000 $\mu$M and 41 linearly spaced from 450 to 650 $\mu$M. Twenty-five logarithmic capacities from $10^{-4}$ to 0.8, two direct-route settings, and $J=0,0.4$ give 8,200 paired comparisons. Totals are $P_T=5$, $S_T=10$, $D_T=2.5\epsilon_P$, and $E_T=1.25\epsilon_P$. Random verification additionally samples 60 cases with log-uniform $G\in[10,10^4]$, $P_T\in[10^{-0.3},10]$, $S_T\in[1,100]$, $D_T,E_T\in[10^{-3},10^{-0.3}]$, and uniform $\lambda\in[1,6]$, $\sigma\in[1,5]$, $J\in[0,0.4]$. Each case uses two free-pool initial guesses. Independent BDF integrations start with free unmodified PII and GS, unliganded GlnD and GlnE, and all complexes absent. They use $G=530$, $P_T=5$, $S_T=10$, $D_T=0.5$, $E_T=0.25$ and $(\lambda,\sigma,J)=(1,1,0),(4,4,0),(4,4,0.1)$, with relative tolerance $2\times10^{-8}$ and absolute tolerance $2\times10^{-11}$. Time units are arbitrary, not biological seconds. Numerical maxima and tolerances are supplied in the machine-readable summary. These integrations and root comparisons do not establish global stability or uniqueness.

\section*{A11 Cascade-wide bounds and observation boundaries}

\textbf{Total PII distributions.} Suppose two mechanisms have the same normalized free PII distribution $\pi$ at fixed free glutamine and identical protein totals. Let $b_m$ be the fraction of PII bound to any enzyme in model $m$, and $q_m$ its bound-state distribution. Then $p_m=(1-b_m)\pi+b_mq_m$. Since a GlnD contains at most one PII and a GlnE at most two, including GlnE--GS complexes,
\[
 b_m\leq \epsilon_P=\min\{1,(D_T+2E_T)/P_T\}.
\]
The two distributions share at least $1-\max(b_1,b_2)$ mass in $\pi$, so their total-variation distance, and hence their mean modified-fraction difference, is at most $\epsilon_P$. This extends the one-target sequestration argument to a downstream enzyme that can bind two upstream trimers.

\textbf{GS propagation.} For $\epsilon_P<1$, the free PII pool lies in $T\in[P_T(1-\epsilon_P),P_T]$. At common free-state fractions write $p=cT$ and $u=dT$. The specified GlnE drive is
\[
 r(T)=k\frac{cT+\beta g+\beta'cgT/\alpha_1}{dT}
     =a+\frac{b}{T},\quad
 a=\frac{kc(1+\beta'g/\alpha_1)}d,\quad b=\frac{k\beta g}d.
\]
The diagonal transformations preserve $a,b$. Positivity gives $|\Delta\log r|\leq L=\log[1/(1-\epsilon_P)]$. For a shared positive GS ladder, $d\,\operatorname{logit}\theta/d\log r=\nu(r)\leq\nu_*$, where $\nu_*=1$ for independent sites and 12 suffices for any positive twelve-site ladder. Integrating and maximizing the difference of two logistic values at fixed logit separation gives
\[
 |\Delta\theta_{GS,\mathrm{free}}|\leq\tanh(\nu_*L/4).
\]
Bound GS is at most $\epsilon_S=\min(1,E_T/S_T)$ of total GS. Coupling the two free-state contributions, with the remaining bound mass allowed to vary arbitrarily, adds at most $\epsilon_S$. This yields the main-text bound. If $\epsilon_P\geq1$, only the trivial unit bound is guaranteed by this argument. If $\beta=0$, $b=0$ and free GS distributions coincide exactly, leaving only $\epsilon_S$. A tighter condition-specific bound replaces the hyperbolic tangent by the difference between the GS fractions evaluated at the two allowed free-pool endpoints. The sweep verifies both bounds. No derivative or fitted Hill-coefficient bound follows from these fraction bounds.

At $\epsilon_P=0.01$ and $\epsilon_S=0.001$, the conservative bounds are approximately 0.003513 for independent GS and 0.03114 for the unrestricted twelve-site ladder. For a tolerated fraction difference $\delta>\epsilon_S$, inversion yields the sufficient protein-pool requirement
\[
 \epsilon_P\leq1-\exp\left[-\frac{4}{\nu_*}\operatorname{artanh}(\delta-\epsilon_S)\right].
\]
With $\delta=0.01$ and $\epsilon_S=0.001$, this gives 0.03536 for independent sites and 0.002996 for the general twelve-site ladder. These are allowed upper limits, not expected separations or empirical noise estimates. A null total-readout result within those limits cannot select an equivalence member. Conversely, a difference exceeding the applicable bound rejects at least one of the shared-topology, input-control, pool-size, or equivalence assumptions; it does not identify which one.

\textbf{Recognition and productive-route checks.} Replacing $i/3$ by $(i/3)^\chi$, $\chi=0.5,1,2$, preserves $u=dT$ and hence the proof. Ninety selected combinations verify this extension. In contrast, making EPU catalytically productive adds $k\eta pu/\alpha_2$ to the ATase numerator. After division by $u$, the drive gains $k\eta p/\alpha_2$, destroying the silent-state $\alpha_2$ invariance and the particular $a+b/T$ closure. For $\alpha_2=2.17$ versus 8.85, 21 values of $\eta$ (zero and twenty logarithmic values from $10^{-4}$ to one) and 41 glutamine doses give 861 comparisons. The largest selected fixed-free-pool GS fraction separation is 0.15449 at $G=447.214\,\mu$M and $\eta=1$. The corresponding total-GS difference is 0.15301 with $P_T=5,S_T=10,D_T=0.05,E_T=0.025$. For the original $\lambda/\sigma$ pair at common $\alpha_2$, however, $a,b,c$ are unchanged and $r=a+b/T+cT$ satisfies $|d\log r/d\log T|=|cT-b/T|/(a+b/T+cT)\leq1$. Thus the general fraction bound survives this extension; only the special $b=0$ exact-invariance conclusion is lost when $c>0$. An additional 360 paired conditions verify this surviving bound. The productive route is a countermodel, not an experimentally demonstrated activity of EPU. The earlier bound-complex exchange construction in A8 supplies a different topology boundary.

\textbf{Conserved glutamine.} If only regulatory glutamine binding is added to input conservation, each GlnD binds at most two glutamines and each GlnE at most one, whether or not a catalytic target is attached. Thus
\[
 \max(0,G_T-2D_T-E_T)\leq G\leq G_T.
\]
This is not a slope bound. The implementation solves the buffered-input protein balances inside a scalar glutamine-balance root. Across 31 total inputs from 10 to 2000 $\mu$M at each of two protein scales, the two transformed models acquire different free inputs. The larger-scale scenario $P_T=50,S_T=100,D_T=5,E_T=2.5$ gives a selected maximum free-PII total-variation difference $6.34\times10^{-4}$, illustrating a small but nonzero loss of the exact fixed-free-input identity. This steady-state extension has not been independently validated by integration of a glutamine-dynamic reaction system. Metabolic glutamine consumption, transport, and other binding partners are outside its scope.

\section*{A12 A resolution limit for complementary stationary measurements}

A stationary occupancy assay distinguishes the specific diagonal families only within its assumed state space. To see the limit, split a regulatory state of weight $w$ into two unobserved conformations with the same ligand stoichiometry and weights $fw$ and $(1-f)w$, $0<f<1$. Give the first conformation opposing specificities $a^{\pm}/f$ and the second no catalytic branches. The summed binding weight remains $w$, and the opposing catalytic numerators remain $wa^{\pm}$. Thus ligand occupancy and the target stationary response are both preserved in this reduced construction. If target-complex association coefficients on the active conformation are also $h^{\pm}/f$, their stationary contributions are preserved as well. Positive active-branch rates can be realized as in A9; the inactive conformation has absent catalytic edges rather than a negative rate. A connected reversible conformational network can realize the prescribed free-state weights.

This is a counterexample, not evidence for hidden GlnD/GlnE conformations or identical transients. Fig 4 distinguishes only its specified classes and nuisance ranges. A larger state space requires another discriminating observable or independent restrictions on catalytic states.


\section*{A13 Paired-paralog reconstruction and deterministic uncertainty}

At the subunit level, let $N_K$ count all GlnK subunits and $N_{K,b}$ those bound to AmtB. If all bound GlnK subunits are unmodified, $N_{K,b}\leq N_K(1-q_K)$, irrespective of their assembly into homo- or heterooligomers. Thus $s_K=N_{K,b}/N_K\leq1-q_K$. For fixed homotrimers, this protein fraction equals the particle fraction used below. In mixed particles, the peptide-based ceiling does not by itself specify the number of bound particles or the distribution of their compositions. The paired reconstruction below retains the homotrimer and free-GlnB assumptions; the observed occurrence of mixed trimers (main reference to van Heeswijk et al., 2000) makes these assay controls substantive.

Before specifying kinetics, let $p_i$ be any distribution of the four trimer modification counts and $q=\sum_i ip_i/3$. Since $q\leq1-p_0$, binding restricted to count zero gives $s\leq p_0\leq1-q$. The last bound is sharp: a distribution with mass $1-q$ at count zero and mass $q$ at count three attains it, without any independence or stationarity assumption. This does not prove that all unmodified trimers are actually bound. The source observations constrain $q$ only through their stated peptide calibration model.

The normalization and uncertainty analysis for the independent peptide measurements is given in S2 Appendix, D6.


Let the free GlnB ladder have binomial weights $\binom3i u^i$ and free GlnK have weights $\binom3i(\rho u)^i$. Only unmodified GlnK binds AmtB, and the bound class cannot be modified. At steady state the terminal binding branch has zero net flux; cutting each modification edge then gives the free binomial ladder. This result allows finite association and dissociation rates but assumes no bound catalytic route. A buffered free AmtB construction has five GlnK weights
\[
 (1,3v,3v^2,v^3,\beta),\qquad v=\rho u,\quad\beta\geq0.
\]
Their normalizer is $(1+v)^3+\beta$. The bound fraction is $s=\beta/[(1+v)^3+\beta]$, while total modified-site fraction is $q_K=(1-s)v/(1+v)$. Since $q_B=u/(1+u)$, elimination of $u$ gives the main-text reconstruction. The algebra does not require determining $\beta$ or the GlnD switching law. A conserved finite AmtB pool instead determines $\beta$ through free transporter; the reconstruction still holds at a stationary point satisfying these same branch and ladder conditions.

Without binding, the control obeys $\rho=\operatorname{odds}(q_K^0)/\operatorname{odds}(q_B^0)$. The control may have another driving ratio $u$ but must preserve relative catalytic specificity and ladder shape. Interior intervals $q_B^0\in[B_-,B_+]$, $q_K^0\in[K_-,K_+]$ therefore give
\[
 \rho_- =\frac{\operatorname{odds}(K_-)}{\operatorname{odds}(B_+)},\qquad
 \rho_+ =\frac{\operatorname{odds}(K_+)}{\operatorname{odds}(B_-)}.
\]
Endpoint saturation makes odds calibration poorly conditioned; it is not repaired by assuming exact endpoints from noisy data.

For assay intervals $q_B\in[b_-,b_+]$, $q_K\in[k_-,k_+]$, impose an independently justified model-error bound
\[
 |q_{K,\rm free}-F_\rho(q_B)|\leq\delta.
\]
Monotonicity yields $f_- =\max(0,F_{\rho_-}(b_-)-\delta)$ and $f_+=\min(1,F_{\rho_+}(b_+)+\delta)$. If $k_->f_+$ the assumptions and input intervals are jointly incompatible. Otherwise, for $f_+>0$,
\[
 s_-=\begin{cases}0,&f_-\leq k_+,\\1-k_+/f_-,&f_->k_+,\end{cases}
 \qquad s_+=1-k_-/f_+.
\]
If $f_+=0$ and the total interval includes zero, every $s\in[0,1]$ is feasible. The endpoints are attainable within the stipulated independent interval box; correlations or shared calibration biases can shrink the feasible set and require a different propagation. Confidence coverage follows only if the joint input box and model-error bound have the claimed coverage. A measurement SE alone does not justify $\delta$.

Two thousand random stationary constructions are checked by independent five-state generator solves. Ten thousand randomly generated admissible configurations check interval containment, including separate zero and incompatible cases. Neither verification establishes physiological stationarity. The source data are transient paired-peptide measurements, and native binding also depends on ATP and 2-oxoglutarate. A practical validation would compare the reconstructed fraction with a direct AmtB--GlnK binding assay under controlled inputs, avoiding saturation in the calibration and testing that both reporter fractions have stabilized.


\section*{A14 Calibrated curves and sharp bounds beyond binomial ladders}
Let $p_i(u)=a_i u^i/\sum_{j=0}^n a_j u^j$, with all $a_i>0$, and $m(u)=\sum_i ip_i(u)/n$. The established variance identity gives $d m/d\log u=\operatorname{Var}(i)/n>0$ for finite positive $u$. Two free ladders sharing one scalar drive therefore define an increasing curve $C=m_K\circ m_B^{-1}$. A no-binding titration can calibrate this curve without identifying the ladders separately. Independent sites and a fixed odds ratio are sufficient special cases, not necessary conditions. Transfer of the calibrated curve across assay conditions remains an experimental requirement; multiple independently varying inputs need not collapse to one curve.

Suppose total reporter modification lies in $[B_L,B_U]$, at most a fraction $\epsilon<1$ of reporter protein is sequestered, and bound-reporter modification is otherwise unrestricted. Conservation gives the sharp interval
\[
q_{B,\mathrm{free}}\in\left[\max\left(0,\frac{B_L-\epsilon}{1-\epsilon}\right),\min\left(1,\frac{B_U}{1-\epsilon}\right)\right].
\]
For $\epsilon=1$, the free fraction is unrestricted. Passing this interval through a calibrated curve envelope gives an allowed free-target interval $[F_L,F_U]$. The envelope must cover both calibration error and biological transfer error. Endpoint evaluation is valid for monotone envelopes; otherwise their extrema over the reporter interval are required.

Let bound-target modification lie in $[H_L,H_U]$ and total-target modification in $[Q_L,Q_U]$. For bound protein fraction $s$, conservation gives $q=(1-s)f+s h$. Its sharp identified set is $[0,1]$ intersected with
\[
(H_L-F_L)s\leq Q_U-F_L,\qquad
(F_U-H_U)s\leq F_U-Q_L.
\]
Indeed, at fixed $s$, attainable total modification fills exactly $[(1-s)F_L+sH_L,(1-s)F_U+sH_U]$. This interval overlaps the observed interval precisely when the displayed inequalities hold. Every feasible $s$ is attainable within the rectangular input set. Dependencies between input uncertainties can tighten the set. Empty intersection rejects the joint assumptions, without identifying which assumption failed. Equal free and bound modification eliminates the contrast needed to identify $s$.

The mixture statement is instantaneous and applies to protein subunits irrespective of oligomer composition. Stationarity or controlled composition may still be required to justify the supplied calibration. Thus the result removes these assumptions from the conservation step, not from every possible biological application.

For synthetic $f\in[0.70,0.80]$ and $q\in[0.35,0.39]$, $h=0$ gives $s\in[0.443,0.563]$; permitting $h\in[0,0.15]$ gives $[0.443,0.692]$. With unrestricted $h\in[0,1]$, only the upper endpoint becomes one. These values illustrate partial identification, not measured occupancy. Three thousand independent linear programs verify mixture interval endpoints, 500 verify reporter intervals, and 1,000 arbitrary positive ladder pairs verify curve inversion. The elementary mixture argument and variance identity are established mathematical ingredients; the extension here specifies which calibration and partition measurements make the proposed assay informative.

\section*{A15 Established theory and present contributions}

\begin{center}\small
\begin{tabular}{>{\raggedright\arraybackslash}p{0.20\textwidth}>{\raggedright\arraybackslash}p{0.32\textwidth}>{\raggedright\arraybackslash}p{0.38\textwidth}}
\toprule
Result & Established foundation & Contribution and restriction here\\\midrule
Stationary ladder & Rational parameterization and catalytic-complex elimination: Thomson and Gunawardena (2009); Feliu and Wiuf (2013). & A reduction under specified shared-enzyme and regulatory-exchange restrictions. It is a method, not a new general parameterization theorem.\\[0.7em]
Adjacent-state contrast & Gunawardena (2007), Eqs. 3--5, doi:10.1529/biophysj.107.110866. & Application to fixed-ladder and refractory alternatives. The invariant itself is established and its constancy does not prove a native mechanism.\\[0.7em]
Variance and sensitivity & Binding-polynomial differentiation; Adair (1925), Briggs (1984), Abeliovich (2005). & Separating physical local slopes from fitted or range coefficients prevents estimator-dependent attribution errors.\\[0.7em]
Regulatory compensation & Scaling non-identifiability: Castro and de Boer (2020), with limits clarified by Villaverde and Massonis (2021). & Explicit GlnD and source-topology GlnE families preserve target output but separate binding and kinetics. They do not enumerate all equivalent mechanisms.\\[0.7em]
Cascade fraction bound & Conservation, bounded support and logit sensitivity inequalities. & Combining stoichiometric sequestration capacity with downstream free-pool elasticity gives a sufficient assay concentration criterion in the specified closed cascade. It is not a universal Hill-slope bound.\\[0.7em]
Measurement design & Design-dependent model discrimination: Silk et al. (2014). & Nuisance-profiled occupancy and binding designs, including depletion correction and overlapping-range limits, for the explicit biochemical families.\\[0.7em]
Paired-paralog assay & Mixture conservation, monotone calibration and elementary odds cancellation; GlnB/GlnK observations of Gosztolai et al. (2017); AmtB binding of Durand and Merrick (2006). & A binding-disabled curve calibration and independently bounded reporter partition and bound-target modification give sharp protein-fraction intervals. The binomial odds relation is a special case; transferable calibration remains required.\\
\bottomrule
\end{tabular}
\end{center}

The cited papers' full bibliographic entries are in the main manuscript. This comparison distinguishes established mathematical ingredients from the present biochemical construction; it does not assert priority over every possible assay using an internal reporter.


\end{document}
